% 自動生成（tools/mkchecklist.py p1）。各章の「ふりかえり」の項目を集めたもの（第13回）
\begingroup\footnotesize
\def\bkcb{\tikz[baseline=-0.5ex]\draw[\bkTF, line width=0.7pt, rounded corners=0.5pt] (0,-0.3em) rectangle (0.8em,0.5em);}
\begin{longtblr}{colspec={Q[c,m,wd=11mm] X[l,m] Q[c,m,wd=10mm] Q[c,m,wd=9mm] Q[c,m,wd=9mm]}, width=\linewidth, rowhead=1, colsep=3pt,
  row{1}={bg=bkVol!14, font=\sffamily\bfseries\scriptsize}, hlines={0.4pt,bkLine}, rowsep=0.7pt}
章 & できるようになったこと & ページ & 1回め & 2回め \\
{\sffamily\bfseries\color{\bkTT}力-0} & クレーンにかかる力を、「動く・回る・ゆがむ・流れる」の4つの顔に分けて言える。 & \pageref{ph:力-0:sec7} & \bkcb & \bkcb \\
 & 何を無視したかを書いてから、ワイヤーが荷を引く力を、7つの段で求められる。 & \pageref{ph:力-0:sec5} & \bkcb & \bkcb \\
 & 身近な「なぜ？」の答えが、この巻のどの章にあるかを探せる。 & \pageref{ph:力-0:sec3} & \bkcb & \bkcb \\
{\sffamily\bfseries\color{\bkTT}力-1} & 位置・速度・加速度をベクトルで書き、微分と積分でつなげられる。 & \pageref{ph:力-1:sec2} & \bkcb & \bkcb \\
 & 加速度が一定の運動の式を導き、ブレーキや投げ上げの計算に使える。 & \pageref{ph:力-1:sec3} & \bkcb & \bkcb \\
 & 投げた物の道すじや、曲がる物の加速度を、成分に分けて求められる。 & \pageref{ph:力-1:sec5} & \bkcb & \bkcb \\
{\sffamily\bfseries\color{\bkTT}力-2} & 物にはたらく力を見つけて、力の図（自由体図）をかける。 & \pageref{ph:力-2:sec1} & \bkcb & \bkcb \\
 & 力の図から運動方程式を立てて、加速度や張力を求められる。 & \pageref{ph:力-2:sec5} & \bkcb & \bkcb \\
 & 摩擦・円運動・空気の抵抗がある運動を、式と数値計算で解ける。 & \pageref{ph:力-2:sec7} & \bkcb & \bkcb \\
{\sffamily\bfseries\color{\bkTT}力-3} & 力積と運動量から、ぶつかるときの平均の力を求められる。 & \pageref{ph:力-3:sec1} & \bkcb & \bkcb \\
 & 運動量保存と反発係数から、衝突のあとの速さを求められる。 & \pageref{ph:力-3:sec4} & \bkcb & \bkcb \\
 & 質量が変わるロケットの速さを求められる。 & \pageref{ph:力-3:sec6} & \bkcb & \bkcb \\
{\sffamily\bfseries\color{\bkTT}力-4} & 高さと速さの関係を、仕事と運動エネルギーで説明できる。 & \pageref{ph:力-4:sec3} & \bkcb & \bkcb \\
 & エネルギーの保存と摩擦の損失から、速さ・高さ・仕事率を求められる。 & \pageref{ph:力-4:sec7} & \bkcb & \bkcb \\
 & $\pxev F=-\mathrm{grad}\,U$ から、位置エネルギーの式を導ける。 & \pageref{ph:力-4:sec6} & \bkcb & \bkcb \\
{\sffamily\bfseries\color{\bkTT}力-5} & モーメントで、てこ・滑車・輪軸の力を計算できる。 & \pageref{ph:力-5:sec1} & \bkcb & \bkcb \\
 & つり合いの式で、反力や倒れる条件を求められる。 & \pageref{ph:力-5:sec6} & \bkcb & \bkcb \\
 & トラスの部材の力を求め、静定か見分けられる。 & \pageref{ph:力-5:sec8} & \bkcb & \bkcb \\
{\sffamily\bfseries\color{\bkTT}力-6} & 細かく分けてたし、慣性モーメントを表計算で求められる。 & \pageref{ph:力-6:sec3} & \bkcb & \bkcb \\
 & 回転の運動方程式と角運動量の保存を使える。 & \pageref{ph:力-6:sec6} & \bkcb & \bkcb \\
 & 角運動量の保存から、ケプラーの法則を導ける。 & \pageref{ph:力-6:sec10} & \bkcb & \bkcb \\
{\sffamily\bfseries\color{\bkTT}力-7} & 転がる車輪の点の速さを求め、$v=r\omega$ を使える。 & \pageref{ph:力-7:sec1} & \bkcb & \bkcb \\
 & 腕の角と手の位置、歯車の速さとトルクを計算できる。 & \pageref{ph:力-7:sec3} & \bkcb & \bkcb \\
 & 転がる加速度や遠心力・コリオリの力を式から導ける。 & \pageref{ph:力-7:sec5b} & \bkcb & \bkcb \\
{\sffamily\bfseries\color{\bkTT}力-8} & 応力とひずみを計算して、棒がどれだけ伸びるか（縮むか）を求められる。 & \pageref{ph:力-8:sec2} & \bkcb & \bkcb \\
 & ばねや棒を直列・並列につないだときのばね定数を求められる。 & \pageref{ph:力-8:series} & \bkcb & \bkcb \\
 & 応力-ひずみ線図を読み、許容応力と安全率から、棒の太さを決められる。 & \pageref{ph:力-8:sec5} & \bkcb & \bkcb \\
 & 原子のばねから「棒もばね」（$k=EA/L$）を導き、熱応力・せん断・ポアソン比を使える。 & \pageref{ph:力-8:sec3} & \bkcb & \bkcb \\
{\sffamily\bfseries\color{\bkTT}力-9} & 引っぱった棒のななめの切り口にかかる、垂直応力とせん断応力を計算できる。 & \pageref{ph:力-9:sec3} & \bkcb & \bkcb \\
 & モールの円をかいて、主応力と最大せん断応力を求められる。 & \pageref{ph:力-9:sec8} & \bkcb & \bkcb \\
 & 応力を行列で表して主応力を固有値として求め、ミーゼス応力で安全率を確かめられる。 & \pageref{ph:力-9:sec7} & \bkcb & \bkcb \\
{\sffamily\bfseries\color{\bkTT}力-10} & 板のたわみが、長さ・厚さ・はば・置く物の重さで、どう変わるかを言える。 & \pageref{ph:力-10:sec1} & \bkcb & \bkcb \\
 & はりを区切って、せん断力・曲げモーメント・たわみを、表計算の足し算で求められる。 & \pageref{ph:力-10:sec4} & \bkcb & \bkcb \\
 & つり合いと曲げの式から $\sigma=My/I$ とたわみの微分方程式を導き、代表的なはりのたわみを計算できる。 & \pageref{ph:力-10:sec8} & \bkcb & \bkcb \\
{\sffamily\bfseries\color{\bkTT}力-11} & 棒を押す・ねじる・内からふくらませると、どこがどう変わるかを、実験と表から言える。 & \pageref{ph:力-11:sec1} & \bkcb & \bkcb \\
 & 軸が伝える動力とねじりモーメント、圧力容器の壁の応力、柱が曲がり始める力を計算できる。 & \pageref{ph:力-11:sec5} & \bkcb & \bkcb \\
 & ねじり・ひずみエネルギー・座屈の式を導き、カスティリアーノの定理でたわみを求められる。 & \pageref{ph:力-11:sec9} & \bkcb & \bkcb \\
{\sffamily\bfseries\color{\bkTT}力-12} & 圧力を計算し、水圧 $p=p_0+\rho gh$ を（ゲージ圧と絶対圧を区別して）使える。 & \pageref{ph:力-12:sec3} & \bkcb & \bkcb \\
 & 浮力が押しのけた水の重さになるわけを説明し、浮き沈みを計算できる。 & \pageref{ph:力-12:sec7} & \bkcb & \bkcb \\
 & 油圧・壁や水門にかかる力（圧力の中心）・船の安定を計算できる。 & \pageref{ph:力-12:sec6} & \bkcb & \bkcb \\
{\sffamily\bfseries\color{\bkTT}力-13} & 流量 $Q=Av$ と連続の式で、管や蛇口の水の速さと太さを求められる。 & \pageref{ph:力-13:sec1} & \bkcb & \bkcb \\
 & ベルヌーイの式で、速さ・圧力・高さの関係を計算できる（流量計の読みも）。 & \pageref{ph:力-13:sec2} & \bkcb & \bkcb \\
 & バランス式（運動量の法則）で、噴流・曲がり管・ポンプの力を計算できる。 & \pageref{ph:力-13:sec8} & \bkcb & \bkcb \\
{\sffamily\bfseries\color{\bkTT}力-14} & レイノルズ数を計算して、管の流れが層流か乱流かを見分けられる。 & \pageref{ph:力-14:sec7} & \bkcb & \bkcb \\
 & 管摩擦係数と損失係数から、管の圧力損失とポンプの動力を求められる。 & \pageref{ph:力-14:sec9} & \bkcb & \bkcb \\
 & 流体の力のつり合いから、ハーゲン・ポアズイユの式を導ける。 & \pageref{ph:力-14:sec8} & \bkcb & \bkcb \\
{\sffamily\bfseries\color{\bkTT}力-15} & 単位（次元）をそろえて、物理の量の関係の形を見つけられる。 & \pageref{ph:力-15:sec4} & \bkcb & \bkcb \\
 & 抗力係数と揚力係数を使って、車・ボール・飛行機にはたらく力を見積もれる。 & \pageref{ph:力-15:sec5} & \bkcb & \bkcb \\
 & 無次元数をそろえて、模型の実験から本物にはたらく力を求められる。 & \pageref{ph:力-15:sec7} & \bkcb & \bkcb \\
{\sffamily\bfseries\color{\bkTT}力-16} & 形を要素に分け、ばねをつないだものとして、行列の式 $K\vec u=\vec F$ で解く考え方がわかる。 & \pageref{ph:力-16:sec2} & \bkcb & \bkcb \\
 & 時間を細かく刻んで衝突を計算する考え方と、刻みの大きさの決め方がわかる。 & \pageref{ph:力-16:sec3} & \bkcb & \bkcb \\
 & 要素を細かくしたり、手計算とくらべたりして、コンピュータの答えを確かめられる。 & \pageref{ph:力-16:sec5} & \bkcb & \bkcb \\
\end{longtblr}
\endgroup
